% Luc Destombe --- licence CC BY-NC-SA 4.0
% https://creativecommons.org/licenses/by-nc-sa/4.0/deed.fr
% Fichier autonome : compiler avec pdflatex, deux fois (signets, nombre de pages).
\documentclass[11pt,a4paper]{article}
\makeatletter
\newif\iflycee@unecolonne
\lycee@unecolonnetrue
\newif\iflycee@cm
\newif\iflycee@palatino

\RequirePackage[utf8]{inputenc}
\RequirePackage[T1]{fontenc}
\RequirePackage[french]{babel}
\RequirePackage{etoolbox}

\newcommand{\ShorthandsOff}{\@ifundefined{shorthandoff}{}{\shorthandoff{;:!?}}}
\newcommand{\ShorthandsOn}{\@ifundefined{shorthandon}{}{\shorthandon{;:!?}}}
\ShorthandsOff
\AtBeginDocument{\ShorthandsOn}
\AtBeginEnvironment{tikzpicture}{\ShorthandsOff}

\RequirePackage{geometry}
\RequirePackage{fancyhdr}
\RequirePackage{lastpage}
\RequirePackage{multicol}
\RequirePackage{array}
\RequirePackage{enumitem}
\RequirePackage{xcolor}
\RequirePackage{needspace}

\RequirePackage{amsmath,amssymb}
\iflycee@cm\else\RequirePackage{fourier}\fi
\DeclareMathAlphabet{\mathcal}{OMS}{cmsy}{m}{n}
\renewcommand{\ttdefault}{lmtt}
\begingroup\catcode`\_=\active \catcode`\^=\active
\global\let_\sb \global\let^\sp
\endgroup
\AtBeginDocument{\catcode`\_=12 \mathcode`\_="8000
  \catcode`\^=12 \mathcode`\^="8000 }
\newlength{\retraitcalculs}
\setlength{\retraitcalculs}{2em}

\RequirePackage{tikz}
\usetikzlibrary{arrows.meta,calc,babel,shapes.misc}
\RequirePackage[most]{tcolorbox}
\tcbuselibrary{skins,breakable}

\definecolor{coulDef}{HTML}{1F4E79}
\definecolor{coulProp}{HTML}{7030A0}
\definecolor{coulRem}{HTML}{C55A11}
\definecolor{coulEx}{HTML}{2E7D32}
\definecolor{coulExo}{HTML}{B03A2E}
\definecolor{coulPart}{HTML}{1F4E79}
\definecolor{coulMeth}{HTML}{1B6E4F}
\definecolor{coulCourbe}{HTML}{0B5ED7}
\definecolor{coulCourbeB}{HTML}{C00000}
\definecolor{coulCourbeC}{HTML}{2E7D32}
\definecolor{coulGrille}{HTML}{8C8C8C}
\definecolor{coulRep}{HTML}{1B6E4F}
\definecolor{coulBar}{HTML}{2E7D32}

\newcommand{\R}{\mathbb{R}}
\newcommand{\Rs}{\mathbb{R}^{*}}
\newcommand{\N}{\mathbb{N}}
\newcommand{\Ns}{\mathbb{N}^{*}}
\newcommand{\Z}{\mathbb{Z}}
\newcommand{\Q}{\mathbb{Q}}
\renewcommand{\le}{\leqslant}
\renewcommand{\ge}{\geqslant}
\renewcommand{\approx}{\simeq}
\newcommand{\vect}[1]{\overrightarrow{#1}}
\newcommand{\norme}[1]{\left\lVert #1 \right\rVert}
\newcommand{\intervalle}[2]{\left[#1\,;#2\right]}
\RequirePackage{cancel}
\renewcommand{\CancelColor}{\color{coulExo}}
\newcommand{\fois}[1]{\times{\color{coulExo}#1}}
\newcommand{\barre}[1]{\cancel{#1}}
\newlength{\lycee@eqh}\newlength{\lycee@eqd}
\newsavebox{\lycee@eqa}\newsavebox{\lycee@eqb}\newsavebox{\lycee@eqc}
\newcommand{\lycee@eqcol}[1]{\hbox{$\renewcommand{\arraystretch}{1.35}%
  \begin{array}[t]{@{}>{\displaystyle}r@{}>{\displaystyle{}}l@{}}#1\end{array}$}}
\newcommand{\lycee@eqmesure}[1]{%
  \ifdim\ht#1>\lycee@eqh\setlength{\lycee@eqh}{\ht#1}\fi
  \ifdim\dp#1>\lycee@eqd\setlength{\lycee@eqd}{\dp#1}\fi}
\newcommand{\lycee@eqtrait}{\hspace{1.2em}\tikz[baseline=0pt]\draw[dash pattern=on 2pt off 2pt,
  line width=0.5pt] (0,\lycee@eqh) -- (0,-\lycee@eqd);\hspace{1.2em}}
\newcommand{\eqproduit}[2]{\par\smallskip\noindent C'est une équation produit :
  $A\times B=0 \iff A=0$ ou $B=0$.\par\nopagebreak\smallskip
  \begingroup\setlength{\lycee@eqh}{0pt}\setlength{\lycee@eqd}{0pt}%
  \setbox\lycee@eqa=\lycee@eqcol{#1}\setbox\lycee@eqb=\lycee@eqcol{#2}%
  \lycee@eqmesure\lycee@eqa\lycee@eqmesure\lycee@eqb%
  \noindent\hspace*{\retraitcalculs}\usebox\lycee@eqa\lycee@eqtrait\usebox\lycee@eqb\endgroup\par\smallskip}
\newcommand{\eqproduittrois}[3]{\par\smallskip\noindent C'est une équation produit :
  $A\times B\times C=0 \iff A=0$ ou $B=0$ ou $C=0$.\par\nopagebreak\smallskip
  \begingroup\setlength{\lycee@eqh}{0pt}\setlength{\lycee@eqd}{0pt}%
  \setbox\lycee@eqa=\lycee@eqcol{#1}\setbox\lycee@eqb=\lycee@eqcol{#2}\setbox\lycee@eqc=\lycee@eqcol{#3}%
  \lycee@eqmesure\lycee@eqa\lycee@eqmesure\lycee@eqb\lycee@eqmesure\lycee@eqc%
  \noindent\hspace*{\retraitcalculs}\usebox\lycee@eqa\lycee@eqtrait\usebox\lycee@eqb\lycee@eqtrait\usebox\lycee@eqc\endgroup
  \par\smallskip}

\tikzset{grille/.style={coulGrille,line width=0.4pt}}
\newcommand{\quadrillage}[4]{\draw[grille,step=1] (#1,#2) grid (#3,#4);}
\tikzset{
  croix/.style={cross out,draw,line width=0.6pt,inner sep=0pt,minimum size=4pt},
  croixdroite/.style={croix,rotate=45,line width=0.8pt},
}
\newcommand{\pt}[3]{\path (#1) node[croix]{} node[#2,font=\small,inner sep=2.5pt]{$#3$};}

\newcommand{\lycee@repere}[4]{%
  \draw[grille,step=1] (#1,#2) grid (#3,#4);
  \draw[-{Stealth[length=2mm]},thick] (#1,0)--({#3+0.4},0) node[below left=-1pt]{$x$};
  \draw[-{Stealth[length=2mm]},thick] (0,#2)--(0,{#4+0.4}) node[below left=-1pt]{$y$};
  \node[below left,font=\scriptsize,inner sep=1.5pt] at (0,0) {$0$};}
\newcommand{\lycee@gradx}[1]{\foreach \k/\l in {#1}{%
  \draw (\k,0.07)--(\k,-0.07) node[below,font=\scriptsize,inner sep=1.5pt]{$\l$};}}
\newcommand{\lycee@grady}[1]{\foreach \k/\l in {#1}{%
  \draw (0.07,\k)--(-0.07,\k) node[left,font=\scriptsize,inner sep=1.5pt]{$\l$};}}
\newcommand{\lycee@intff}[2]{\left[#1\,;#2\right]}
\newcommand{\lycee@intfo}[2]{\left[#1\,;#2\right[}
\newcommand{\lycee@intof}[2]{\left]#1\,;#2\right]}
\newcommand{\lycee@intoo}[2]{\left]#1\,;#2\right[}
\newcommand{\lycee@ens}[1]{\left\{#1\right\}}
\AtBeginDocument{%
  \providecommand{\repere}{\lycee@repere}%
  \providecommand{\gradx}{\lycee@gradx}%
  \providecommand{\grady}{\lycee@grady}%
  \providecommand{\Cf}{\mathcal{C}\sb{f}}%
  \providecommand{\Cg}{\mathcal{C}\sb{g}}%
  \providecommand{\intff}{\lycee@intff}%
  \providecommand{\intfo}{\lycee@intfo}%
  \providecommand{\intof}{\lycee@intof}%
  \providecommand{\intoo}{\lycee@intoo}%
  \providecommand{\ens}{\lycee@ens}}

\newlist{questions}{enumerate}{3}
\setlist[questions,1]{label=\arabic*\textdegree),leftmargin=*,itemsep=2pt,topsep=3pt}
\setlist[questions,2]{label=\alph*),leftmargin=*,itemsep=1pt,topsep=2pt}
\setlist[questions,3]{label=\roman*),leftmargin=*,itemsep=1pt,topsep=1pt}
\setlist[itemize]{label=\textbullet,leftmargin=*,itemsep=2pt,topsep=3pt}

\newcommand{\besoinplace}[1]{\par\penalty\@M\begingroup
  \dimen@=#1\relax\dimen@ii=\pagegoal\advance\dimen@ii-\pagetotal
  \ifdim\dimen@>\dimen@ii\ifdim\dimen@ii>\z@\vfil\fi\break\fi\endgroup}

\newcommand{\lycee@pied}{\thepage/\pageref{LastPage}}
\newcommand{\entete}[2]{%
  \pagestyle{fancy}\fancyhf{}%
  \fancyhead[L]{\footnotesize\color{coulPart}\textbf{#1}}%
  \fancyhead[R]{\footnotesize\color{coulPart}\textbf{#2}}%
  \fancyfoot[C]{\footnotesize\lycee@pied}%
  \renewcommand{\headrulewidth}{0.4pt}%
  \renewcommand{\headrule}{\hbox to\headwidth{%
    \color{coulPart!40}\leaders\hrule height \headrulewidth\hfill}}}

\RequirePackage{ccicons}
\newcommand{\lycee@licence}{{\scriptsize\color{gray}%
  \copyright\ Luc Destombe --- \ccbyncsa\ CC BY-NC-SA 4.0}}
\newif\iflycee@licence \lycee@licencetrue
\newcommand{\sanslicence}{\lycee@licencefalse}
\AtBeginDocument{\iflycee@licence\AddToHookNext{shipout/foreground}{%
  \put(0,0){\rlap{\kern\dimexpr1in+\hoffset+\oddsidemargin+\textwidth\relax
    \llap{\raisebox{-\dimexpr1in+\voffset+\topmargin+\headheight+\headsep
      +\textheight+\footskip\relax}{\lycee@licence}}}}}\fi}

\newcommand{\lycee@site}{faire-des-maths.fr}
\newcommand{\lycee@version}{{\scriptsize\color{gray}\lycee@site
  \ifdefined\lyceeversion\ --- Version~\lyceeversion\fi}}
\AtBeginDocument{\iflycee@licence\AddToHookNext{shipout/foreground}{%
  \put(0,0){\rlap{\kern\dimexpr1in+\hoffset+\oddsidemargin\relax
    \raisebox{-\dimexpr1in+\voffset+\topmargin+\headheight+\headsep
      +\textheight+\footskip\relax}{\lycee@version}}}}\fi}

\newcommand{\formulesserrees}{%
  \setlength{\abovedisplayskip}{3pt}\setlength{\belowdisplayskip}{3pt}%
  \setlength{\abovedisplayshortskip}{2pt}\setlength{\belowdisplayshortskip}{3pt}}

\setlength{\parindent}{0pt}

\RequirePackage[implicit=false,hidelinks,pdfencoding=auto,psdextra,bookmarksopen,
  bookmarksopenlevel=1,pdfcreator={},pdfproducer={}]{hyperref}
\RequirePackage{bookmark}
\hypersetup{pdfauthor={Luc Destombe},pdfkeywords={CC BY-NC-SA 4.0}}
\newcounter{lycee@signet}
\newcommand{\signet}[3][]{\stepcounter{lycee@signet}%
  \edef\lycee@dest{\if\relax\detokenize{#1}\relax
    signet.\arabic{lycee@signet}\else#1\fi}%
  \hypertarget{\lycee@dest}{}\bookmark[level=#2,dest=\lycee@dest]{#3}}
\pdfstringdefDisableCommands{%
  \def\R{R}\def\N{N}\def\Z{Z}\def\Q{Q}\def\vect#1{#1}\def\Cf{Cf}\def\Cg{Cg}\def\fois#1{×#1}%
  \def\barre#1{#1}\def\intervalle#1#2{[#1;#2]}\def\intff#1#2{[#1;#2]}%
  \def\textsuperscript#1{#1}\def\dfrac#1#2{#1/#2}\def\frac#1#2{#1/#2}%
  \def\vec#1{#1⃗}}%
\RequirePackage[official]{eurosym}

\tcbset{exercice corrige/.style={
  enhanced,breakable,before upper=\raggedright,
  colback=white,colframe=coulExo,
  boxrule=0.8pt,arc=2pt,
  left=6pt,right=6pt,top=8pt,bottom=5pt,
  before skip=9pt,after skip=4pt,
  coltitle=white,fonttitle=\bfseries\small,
  attach boxed title to top left={xshift=8pt,yshift=-\tcboxedtitleheight/2},
  boxed title style={colback=coulExo,arc=2pt,boxrule=0pt}}}
\newcounter{exo}
\newtcolorbox{exercice}[1][]{exercice corrige,
  phantom={\signet[exercice-\arabic{exo}]{2}{Exercice \arabic{exo}}},
  title={Exercice \theexo\ifx\relax#1\relax\else{} --- #1\fi}}
\newenvironment{exo}[1][\relax]{\refstepcounter{exo}\begin{exercice}[#1]}{\end{exercice}}

\geometry{top=1.8cm,bottom=1cm,left=1cm,right=1cm,headheight=14pt,footskip=0.6cm}
\renewcommand{\lycee@pied}{\thepage}

\renewcommand{\theexo}{\ifnum\value{exo}<10 0\fi\arabic{exo}}
\tcbset{exercice corrige/.style={
  enhanced,breakable,before upper=\raggedright,
  colback=white,colframe=coulExo,
  boxrule=0.9pt,arc=2pt,
  left=8pt,right=8pt,top=8pt,bottom=6pt,
  before skip=8pt,after skip=6pt,
  coltitle=white,fonttitle=\bfseries,
  attach boxed title to top left={xshift=8pt,yshift=-\tcboxedtitleheight/2},
  boxed title style={colback=coulExo,arc=2pt,boxrule=0pt}}}

\newcommand{\entetefiche}[2]{%
  \begin{center}
    \begin{tcolorbox}[enhanced,width=0.92\linewidth,colback=white,colframe=coulPart,
        boxrule=1.4pt,arc=3pt,halign=center,top=6pt,bottom=6pt]
      {\large\bfseries\color{coulPart} CORRIGÉ --- FICHE D'EXERCICES}\\[3pt]
      {\normalsize\bfseries #1}\\[2pt]
      {\normalsize #2}
    \end{tcolorbox}
  \end{center}}

\newcounter{partie}
\newcommand{\partie}[1]{%
  \stepcounter{partie}%
  \besoinplace{8\baselineskip}\vspace{4pt}%
  \begin{tcolorbox}[enhanced,colback=coulPart!8,colframe=coulPart,boxrule=0pt,
      leftrule=3pt,arc=0pt,left=6pt,right=4pt,top=3pt,bottom=3pt,
      before skip=6pt,after skip=2pt]
    \bfseries\color{coulPart}\Roman{partie} --- #1\signet{1}{\Roman{partie} --- #1}
  \end{tcolorbox}}
\newcommand{\partiesansnum}[1]{%
  \besoinplace{8\baselineskip}\vspace{4pt}%
  \begin{tcolorbox}[enhanced,colback=coulPart!8,colframe=coulPart,boxrule=0pt,
      leftrule=3pt,arc=0pt,left=6pt,right=4pt,top=3pt,bottom=3pt,
      before skip=6pt,after skip=2pt]
    \bfseries\color{coulPart}#1\signet{1}{#1}
  \end{tcolorbox}}

\setlist[questions,1]{label=\arabic*\textdegree),leftmargin=*,itemsep=3pt,topsep=3pt}
\setlist[questions,2]{label=\alph*),leftmargin=*,itemsep=2pt,topsep=2pt}
\setlist[questions,3]{label=\roman*),leftmargin=*,itemsep=1pt,topsep=1pt}
\newlist{expr}{enumerate}{1}
\setlist[expr]{label={\Alph*\ $=$},leftmargin=*,itemsep=2pt,topsep=3pt}

\newcommand{\res}[1]{{\color{coulRep}#1}}
\newcommand{\rem}[1]{\par\smallskip{\small\itshape\color{black!60}$\rightarrow$ #1}}
\newenvironment{calculs}{\par\smallskip\noindent$\displaystyle\begin{aligned}[t]}%
  {\end{aligned}$\par\smallskip}
\newcommand{\justif}[1]{\quad\text{\small\itshape\color{black!60}#1}}
\newcommand{\cas}[1]{\par\smallskip\noindent\textbf{\color{coulExo}#1}\vadjust{\nobreak}}
\newcommand{\explique}[2]{\par\smallskip\noindent
  \begin{minipage}[t]{0.6\linewidth}\raggedright #1\end{minipage}\hfill
  \begin{minipage}[t]{0.37\linewidth}\raggedright\leavevmode
    {\small\itshape\color{black!60}#2}\end{minipage}\par}

\setlength{\parskip}{2pt}

\makeatother
\tcbuselibrary{listings}
\entete{Chapitre 11 --- Probabilités --- corrigé}{Seconde}
\setlength{\parindent}{0pt}
\setlength{\jot}{5pt}

\let\casligne\cas
\renewcommand{\cas}[1]{\casligne{#1}\par\nopagebreak}

\newcommand{\Card}{\operatorname{Card}}

\newtcblisting{python}{%
  listing only,enhanced,
  colback=black!3,colframe=black!35,boxrule=0.5pt,arc=1pt,
  left=4pt,right=4pt,top=1pt,bottom=1pt,
  listing options={language=Python,basicstyle=\ttfamily\small,
    keywordstyle=\color{coulMeth}\bfseries,columns=fullflexible,
    keepspaces=true,showstringspaces=false,aboveskip=0pt,belowskip=0pt}}

\tikzset{
  arbre/.style={grow'=right,level distance=2.3cm,
    level 1/.style={sibling distance=2.2cm},
    level 2/.style={sibling distance=1.1cm},
    every node/.style={font=\small}},
  poids/.style={font=\scriptsize,sloped,inner sep=1.5pt},
}

\newcommand{\cotefig}[3][0.36]{\par\noindent
  \begin{minipage}[t]{\dimexpr0.97\linewidth-#1\linewidth\relax}\raggedright #2\end{minipage}\hfill
  \begin{minipage}[t]{#1\linewidth}\vspace{0pt}\centering #3\end{minipage}\par}

\begin{document}

\entetefiche{Chapitre 11 --- Probabilités}{Seconde}

{\small Quand une expérience se fait « au hasard » (dé équilibré, jetons indiscernables,
individu choisi au hasard), on est en situation d'équiprobabilité :
$P(A)=\dfrac{\Card(A)}{\Card(\Omega)}$. On le dit avant de s'en servir.\par}

\partie{Vocabulaire des probabilités}

\begin{exo}[Univers et événements]
\cas{a)}
\explique{$\Omega=\res{\ens{1\,;2\,;3\,;\dots\,;20}}$ et $\Card(\Omega)=\res{20}$.}{les
$20$ numéros possibles}
\cas{b)}
\explique{$A=\res{\ens{6\,;12\,;18}}$.}{$4\times 6=24$ dépasse $20$}
\explique{$B=\res{\ens{17\,;18\,;19\,;20}}$.}{« supérieur ou égal » : $17$ en fait partie}
\cas{c)}
\explique{Aucun numéro de $1$ à $20$ n'est un multiple de $25$, donc
$C=\res{\varnothing}$ : c'est un \res{événement impossible}.}{aucune issue}
\cas{d)}
\explique{$18$ appartient à $A$ et à $B$, mais pas à $C$ : \res{les événements $A$ et
$B$ sont réalisés, $C$ ne l'est pas}.}{on cherche $18$ dans chaque ensemble}
\end{exo}

\begin{exo}[Un univers de mots]
\cas{a)}
\explique{$\Card(\Omega)=\res{8}$.}{on compte les mots}
\cas{b)}
\explique{$A=\res{\ens{\text{mer}\,;\text{rio}\,;\text{ver}\,;\text{mur}\,;\text{roc}}}$.}{les
mots qui contiennent un r}
\explique{$B=\res{\ens{\text{rio}\,;\text{sol}\,;\text{roc}}}$.}{les mots qui contiennent
un o}
\cas{c)}
\explique{Aucun mot ne commence par une voyelle, donc $C=\varnothing$ : « L'événement $C$
est \res{impossible} ».}{a, e, i, o, u, y : aucun mot ne commence ainsi}
\cas{d)}
\explique{Tous les mots ont trois lettres, donc $D=\Omega$ : « L'événement $D$ est
\res{certain} ».}{il contient toutes les issues}
\end{exo}

\begin{exo}[Deux lancers d'une pièce]
\cas{a)}
\explique{$\Omega=\res{\ens{\text{PP}\,;\text{PF}\,;\text{FP}\,;\text{FF}}}$.}{PF et FP
sont deux issues différentes : l'ordre compte}
\cas{b)}
\explique{$A=\res{\ens{\text{PP}\,;\text{PF}\,;\text{FP}}}$.}{« au moins une fois » : une
ou deux fois}
\explique{$B=\res{\ens{\text{PP}\,;\text{FF}}}$.}{deux fois pile ou deux fois face}
\cas{c)}
\explique{FF n'appartient pas à $A$ mais appartient à $B$ : \res{seul l'événement $B$ est
réalisé}.}{FF : aucun pile}
\end{exo}

\partie{Calcul d'une probabilité}

\begin{exo}[Un dé à vingt faces]
\explique{Le dé est équilibré : les $20$ issues ont la même probabilité, et
$\Card(\Omega)=20$.}{on justifie l'équiprobabilité}
\cas{a)}
\explique{$A=\ens{2\,;4\,;\dots\,;20}$ compte $10$ issues.}{un nombre sur deux est pair}
\begin{calculs}
P(A) &= \dfrac{\Card(A)}{\Card(\Omega)} &&\justif{équiprobabilité}\\
P(A) &= \dfrac{10}{20} &&\\
P(A) &= \res{\dfrac{1}{2}} &&\justif{on simplifie par $10$}
\end{calculs}
\cas{b)}
\explique{$B=\ens{6\,;12\,;18}$, donc $P(B)=\res{\dfrac{3}{20}}$.}{$3$ issues sur $20$}
\cas{c)}
\explique{$C=\ens{7\,;14}$, donc $P(C)=\dfrac{2}{20}=\res{\dfrac{1}{10}}$.}{on simplifie
par $2$}
\cas{d)}
\explique{Tous les numéros sont inférieurs ou égaux à $20$, donc $D=\Omega$ et
$P(D)=\res{1}$.}{événement certain}
\end{exo}

\begin{exo}[Des jetons numérotés]
\cas{a)}
\explique{Les jetons sont indiscernables : chacun des $12$ jetons a la même chance d'être
tiré. Les numéros, eux, n'ont pas la même probabilité : il y a plus de jetons
« $1$ » que de jetons « $3$ ». \res{Les issues équiprobables sont donc les $12$
jetons}.}{piège : ne pas compter $3$ issues}
\cas{b)}
\explique{$n$ est impair pour les $5$ jetons « $1$ » et les $3$ jetons « $3$ », soit $8$
jetons.}{$\Card(A)=5+3$}
\begin{calculs}
P(A) &= \dfrac{8}{12} &&\justif{$8$ jetons sur $12$}\\
P(A) &= \dfrac{2\times\barre{4}}{3\times\barre{4}} &&\justif{on simplifie par $4$}\\
P(A) &= \res{\dfrac{2}{3}} &&
\end{calculs}
\explique{$n\ge 2$ pour les $4$ jetons « $2$ » et les $3$ jetons « $3$ », donc
$P(B)=\res{\dfrac{7}{12}}$.}{$\Card(B)=4+3$}
\explique{Aucun jeton ne porte le numéro $4$, donc $P(C)=\res{0}$.}{événement impossible}
\end{exo}

\begin{exo}[Un jeu de 52 cartes]
\explique{La carte est tirée au hasard : équiprobabilité sur les $52$ cartes.}{on justifie
le modèle}
\cas{a)}
\explique{$13$ trèfles : $P=\dfrac{13}{52}=\res{\dfrac{1}{4}}$.}{$52=4\times 13$}
\cas{b)}
\explique{$26$ cartes noires (piques et trèfles) : $P=\dfrac{26}{52}=\res{\dfrac{1}{2}}$.}{$13+13$}
\cas{c)}
\explique{$3$ figures par couleur, soit $12$ figures : $P=\dfrac{12}{52}=\res{\dfrac{3}{13}}$.}{on
simplifie par $4$}
\cas{d)}
\explique{$4$ as : $P=\dfrac{4}{52}=\res{\dfrac{1}{13}}$.}{un as par couleur}
\end{exo}

\begin{exo}[Le sac le plus favorable]
\cas{a)}
\explique{Lina n'a que des billes rouges : $P=\dfrac{4}{4}=\res{1}$.}{événement certain}
\explique{Tom : $6+18=24$ billes, donc $P=\dfrac{6}{24}=\res{\dfrac{1}{4}}$.}{on simplifie
par $6$}
\explique{Inès : $50+2=52$ billes, donc $P=\dfrac{50}{52}=\res{\dfrac{25}{26}}$.}{on
simplifie par $2$}
\explique{$\dfrac{1}{4}<\dfrac{25}{26}<1$ : \res{c'est Lina qui a le plus de chances}
(elle est sûre de tirer une bille rouge).}{$\dfrac{25}{26}\approx 0{,}96$}
\cas{b)}
\explique{Il faut que les $4$ billes rouges de Lina soient le quart du sac, comme pour
Tom.}{$P=\dfrac{1}{4}$}
\explique{Le sac doit contenir $4\times 4=16$ billes : il faut ajouter
$16-4=\res{12}$ billes noires.}{vérification : $\dfrac{4}{16}=\dfrac{1}{4}$}
\end{exo}

\partie{Intersection, réunion, événement contraire}

\begin{exo}[Une roue de loterie]
\cas{a)}
\explique{$A=\ens{4\,;8\,;12\,;16}$ et $B=\ens{13\,;14\,;15\,;16}$.}{on liste d'abord $A$
et $B$}
\explique{$A\cap B=\res{\ens{16}}$.}{seul $16$ est dans les deux}
\explique{$A\cup B=\res{\ens{4\,;8\,;12\,;13\,;14\,;15\,;16}}$.}{$16$ écrit une seule fois}
\explique{$\overline{A}=\res{\ens{1\,;2\,;3\,;5\,;6\,;7\,;9\,;10\,;11\,;13\,;14\,;15}}$.}{les
$12$ numéros qui ne sont pas multiples de $4$}
\explique{$A\cap B\ne\varnothing$, donc \res{$A$ et $B$ ne sont pas
incompatibles}.}{$16$ est dans les deux}
\cas{b)}
\explique{Les secteurs ont la même taille : équiprobabilité, avec $\Card(\Omega)=16$.}{on
justifie le modèle}
\explique{$P(A)=\dfrac{4}{16}=\res{\dfrac{1}{4}}$, $P(B)=\dfrac{4}{16}=\res{\dfrac{1}{4}}$ et
$P(A\cap B)=\res{\dfrac{1}{16}}$.}{$4$, $4$ et $1$ issues}
\cas{c)}
\begin{calculs}
P\big(\overline{A}\big) &= 1-P(A) &&\justif{probabilité du contraire}\\
P\big(\overline{A}\big) &= \dfrac{4}{4}-\dfrac{1}{4} &&\justif{$1=\frac{4}{4}$}\\
P\big(\overline{A}\big) &= \res{\dfrac{3}{4}} &&
\end{calculs}
\begin{calculs}
P(A\cup B) &= P(A)+P(B)-P(A\cap B) &&\justif{propriété 3}\\
P(A\cup B) &= \dfrac{1\fois{4}}{4\fois{4}}+\dfrac{1\fois{4}}{4\fois{4}}-\dfrac{1}{16}
  &&\justif{même dénominateur $16$}\\
P(A\cup B) &= \dfrac{4+4-1}{16} &&\\
P(A\cup B) &= \res{\dfrac{7}{16}} &&
\end{calculs}
\explique{Vérification : $A\cup B$ compte $7$ issues sur $16$, et $\overline{A}$ en compte
$12$, soit $\dfrac{12}{16}=\dfrac{3}{4}$.}{on compte les issues de la question a)}
\end{exo}

\begin{exo}[Traduire par une phrase]
\cas{a)}
\explique{$E\cap A$ : \res{« le spectateur est un étudiant abonné »}.}{« et » : les deux
à la fois}
\explique{$E\cup A$ : \res{« le spectateur est étudiant ou abonné »}, ou les deux.}{« ou » :
au moins l'un des deux}
\explique{$\overline{E}$ : \res{« le spectateur n'est pas étudiant »}.}{le contraire de $E$}
\explique{$E\cap\overline{A}$ : \res{« le spectateur est étudiant et n'a pas
d'abonnement »}.}{$\overline{A}$ : « il n'a pas d'abonnement »}
\cas{b)}
\explique{Léo est étudiant : $E$ est réalisé, donc $E\cup A$ aussi. Il n'a pas
d'abonnement : $A$ et $E\cap A$ ne sont pas réalisés. \res{Les événements $E$ et
$E\cup A$ sont réalisés, $A$ et $E\cap A$ ne le sont pas.}}{pour $E\cup A$, il suffit
d'un des deux}
\end{exo}

\begin{exo}[Défauts de fabrication]
\cas{a)}
\explique{$P(E)=\res{0{,}04}$, $P(B)=\res{0{,}06}$ et $P(E\cap B)=\res{0{,}01}$.}{$4\,\%$,
$6\,\%$ et $1\,\%$ écrits en décimaux}
\cas{b)}
\begin{calculs}
P\big(\overline{E}\big) &= 1-P(E) &&\justif{« pas de défaut d'écran » : contraire de $E$}\\
P\big(\overline{E}\big) &= 1-0{,}04 &&\\
P\big(\overline{E}\big) &= \res{0{,}96} &&
\end{calculs}
\cas{c)}
\begin{calculs}
P(E\cup B) &= P(E)+P(B)-P(E\cap B) &&\justif{« au moins un des deux » : $E\cup B$}\\
P(E\cup B) &= 0{,}04+0{,}06-0{,}01 &&\\
P(E\cup B) &= \res{0{,}09} &&
\end{calculs}
\cas{d)}
\begin{calculs}
P\big(\overline{E\cup B}\big) &= 1-P(E\cup B) &&\justif{« aucun défaut » : contraire de
« au moins un »}\\
P\big(\overline{E\cup B}\big) &= 1-0{,}09 &&\\
P\big(\overline{E\cup B}\big) &= \res{0{,}91} &&
\end{calculs}
\rem{Erreur fréquente : additionner $0{,}04$ et $0{,}06$ sans retirer $0{,}01$ ; les
téléphones qui ont les deux défauts seraient comptés deux fois.}
\end{exo}

\begin{exo}[Des données impossibles]
\cas{a)}
\explique{$P(A)+P\big(\overline{A}\big)$ doit valoir $1$, or $0{,}7+0{,}4=1{,}1$ :
\res{c'est impossible}.}{$P\big(\overline{A}\big)=1-P(A)=0{,}3$}
\cas{b)}
\explique{Si $A$ et $B$ sont incompatibles, $P(A\cup B)=0{,}6+0{,}5=1{,}1$, plus grand que
$1$ : \res{c'est impossible}.}{une probabilité ne dépasse pas $1$}
\cas{c)}
\explique{Les issues de $A\cap B$ sont toutes dans $A$, donc $P(A\cap B)$ ne peut pas
dépasser $P(A)=0{,}3$ : \res{c'est impossible}.}{$A\cap B$ est une partie de $A$}
\cas{d)}
\begin{calculs}
P(A\cup B) &= P(A)+P(B)-P(A\cap B) &&\justif{propriété 3}\\
0{,}7 &= 0{,}5+0{,}35-P(A\cap B) &&\justif{on remplace}\\
0{,}7 &= 0{,}85-P(A\cap B) &&\\
P(A\cap B) &= 0{,}85-0{,}7 &&\justif{on isole $P(A\cap B)$}\\
P(A\cap B) &= \res{0{,}15} &&
\end{calculs}
\end{exo}

\partie{Loi de probabilité}

\begin{exo}[Un dé truqué]
\cas{a)}
\begin{calculs}
P(4) &= 1-(0{,}15+0{,}2+0{,}1+0{,}25+0{,}1) &&\justif{la somme des probabilités vaut $1$}\\
P(4) &= 1-0{,}8 &&\\
P(4) &= \res{0{,}2} &&
\end{calculs}
\cas{b)}
\begin{calculs}
P(A) &= P(2)+P(4)+P(6) &&\justif{propriété 4 : $A=\ens{2\,;4\,;6}$}\\
P(A) &= 0{,}2+0{,}2+0{,}1 &&\\
P(A) &= \res{0{,}5} &&
\end{calculs}
\begin{calculs}
P(B) &= P(4)+P(5)+P(6) &&\justif{« au moins $4$ » : $B=\ens{4\,;5\,;6}$}\\
P(B) &= 0{,}2+0{,}25+0{,}1 &&\\
P(B) &= \res{0{,}55} &&
\end{calculs}
\cas{c)}
\explique{$\overline{B}$ : \res{« obtenir moins de $4$ »}, c'est-à-dire $1$, $2$ ou $3$.}{le
contraire de « au moins $4$ »}
\explique{$P\big(\overline{B}\big)=1-0{,}55=\res{0{,}45}$.}{vérification :
$0{,}15+0{,}2+0{,}1=0{,}45$}
\end{exo}

\begin{exo}[Une roue de quatre couleurs]
\cas{a)}
\begin{calculs}
P(\text{jaune}) &= 1-\dfrac{1}{3}-\dfrac{1}{4}-\dfrac{1}{6} &&\justif{la somme des
probabilités vaut $1$}\\
P(\text{jaune}) &= \dfrac{1\fois{12}}{1\fois{12}}-\dfrac{1\fois{4}}{3\fois{4}}
  -\dfrac{1\fois{3}}{4\fois{3}}-\dfrac{1\fois{2}}{6\fois{2}} &&\justif{même dénominateur $12$}\\
P(\text{jaune}) &= \dfrac{12-4-3-2}{12} &&\\
P(\text{jaune}) &= \dfrac{3}{12} &&\\
P(\text{jaune}) &= \res{\dfrac{1}{4}} &&\justif{on simplifie par $3$}
\end{calculs}
\cas{b)}
\explique{« Ni rouge ni bleue » : vert ou jaune.}{les issues de l'événement}
\begin{calculs}
P(\text{vert ou jaune}) &= \dfrac{1}{4}+\dfrac{1}{4} &&\justif{propriété 4}\\
P(\text{vert ou jaune}) &= \res{\dfrac{1}{2}} &&
\end{calculs}
\cas{c)}
\explique{Les probabilités ne sont pas toutes égales ($\frac{1}{3}\ne\frac{1}{6}$) :
\res{les issues ne sont pas équiprobables}.}{on ne peut pas utiliser
$\dfrac{\Card(A)}{\Card(\Omega)}$}
\end{exo}

\begin{exo}[Attente à la cantine]
\cas{a)}
\begin{calculs}
P(\text{plus de }15) &= 1-(0{,}2+0{,}4+0{,}3) &&\justif{la somme vaut $1$}\\
P(\text{plus de }15) &= \res{0{,}1} &&
\end{calculs}
\cas{b)}
\explique{Moins de $10$ minutes : $0$ à $5$ ou $5$ à $10$, donc
$P=0{,}2+0{,}4=\res{0{,}6}$.}{propriété 4}
\cas{c)}
\explique{Le contraire de « attendre plus de $5$ minutes » est « attendre de $0$ à $5$
minutes », de probabilité $0{,}2$.}{on traduit l'événement contraire}
\explique{$P=1-0{,}2=\res{0{,}8}$.}{vérification : $0{,}4+0{,}3+0{,}1=0{,}8$}
\rem{Cette loi vient de fréquences observées sur un grand nombre de jours : c'est un
modèle, tiré de la loi des grands nombres (partie VIII).}
\end{exo}

\begin{exo}[Une probabilité inconnue]
\cas{a)}
\explique{Les faces $1$ à $5$ ont la probabilité $a$, la face $6$ la probabilité $3a$, et la
somme vaut $1$.}{« trois fois plus grande » : $3a$}
\begin{calculs}
a+a+a+a+a+3a &= 1 &&\justif{somme des probabilités}\\
8a &= 1 &&\justif{on réduit}\\
a &= \res{\dfrac{1}{8}} &&\justif{on divise par $8$}
\end{calculs}
\cas{b)}%
\begin{center}
\renewcommand{\arraystretch}{2.4}
\begin{tabular}{|l|c|c|c|c|c|c|}
\hline
Face & $1$ & $2$ & $3$ & $4$ & $5$ & $6$ \\
\hline
Probabilité & $\dfrac{1}{8}$ & $\dfrac{1}{8}$ & $\dfrac{1}{8}$ & $\dfrac{1}{8}$ &
$\dfrac{1}{8}$ & $\dfrac{3}{8}$ \\
\hline
\end{tabular}
\end{center}
\explique{Vérification : $5\times\dfrac{1}{8}+\dfrac{3}{8}=\dfrac{8}{8}=1$.}{la somme vaut
bien $1$}
\cas{c)}
\begin{calculs}
P(\text{pair}) &= P(2)+P(4)+P(6) &&\justif{propriété 4}\\
P(\text{pair}) &= \dfrac{1}{8}+\dfrac{1}{8}+\dfrac{3}{8} &&\\
P(\text{pair}) &= \res{\dfrac{5}{8}} &&
\end{calculs}
\cas{d)}
\explique{$P(\text{pas }6)=1-\dfrac{3}{8}=\res{\dfrac{5}{8}}$.}{probabilité du contraire}
\end{exo}

\partie{Tableaux croisés d'effectifs}

\begin{exo}[Une salle de concert]
\cas{a)}
\begin{calculs}
\text{gradins} &= 500-300 &&\justif{les autres places}\\
\text{gradins} &= 200 &&
\end{calculs}
\begin{calculs}
\text{fosse, moins de 25 ans} &= 0{,}4\times 300 &&\justif{$40\,\%$ de $300$}\\
\text{fosse, moins de 25 ans} &= 120 &&
\end{calculs}
\begin{calculs}
\text{gradins, moins de 25 ans} &= 0{,}15\times 200 &&\justif{$15\,\%$ de $200$}\\
\text{gradins, moins de 25 ans} &= 30 &&
\end{calculs}
\explique{On complète par soustraction ($300-120=180$, $200-30=170$), puis les totaux des
colonnes par addition.}{chaque ligne a pour somme son total}
\begin{center}
\renewcommand{\arraystretch}{1.5}
\begin{tabular}{|l|c|c|c|}
\hline
 & Moins de 25 ans & 25 ans ou plus & Total \\
\hline
Fosse & $\res{120}$ & $180$ & $300$ \\
\hline
Gradins & $30$ & $170$ & $200$ \\
\hline
Total & $150$ & $350$ & $500$ \\
\hline
\end{tabular}
\end{center}
\cas{b)}
\explique{Le spectateur est choisi au hasard : équiprobabilité sur les $500$
spectateurs.}{propriété 5}
\explique{$P(F)=\dfrac{300}{500}=\res{0{,}6}$, $P(J)=\dfrac{150}{500}=\res{0{,}3}$ et
$P(F\cap J)=\dfrac{120}{500}=\res{0{,}24}$.}{total de ligne, total de colonne, case commune}
\cas{c)}
\begin{calculs}
P(F\cup J) &= P(F)+P(J)-P(F\cap J) &&\justif{propriété 3}\\
P(F\cup J) &= 0{,}6+0{,}3-0{,}24 &&\\
P(F\cup J) &= \res{0{,}66} &&
\end{calculs}
\explique{« En gradins et $25$ ans ou plus » :
$P\big(\overline{F}\cap\overline{J}\big)=\dfrac{170}{500}=\res{0{,}34}$.}{vérification :
$1-0{,}66=0{,}34$}
\end{exo}

\begin{exo}[Musique et dessin]
\cas{a)}
\explique{Musique et dessin : $4$. Musique sans dessin : $14-4=10$. Dessin sans musique :
$10-4=6$. Ni l'un ni l'autre : $32-4-10-6=12$.}{on part de la case « les deux »}
\begin{center}
\renewcommand{\arraystretch}{1.5}
\begin{tabular}{|l|c|c|c|}
\hline
 & Dessin & Pas de dessin & Total \\
\hline
Musique & $\res{4}$ & $10$ & $14$ \\
\hline
Pas de musique & $6$ & $12$ & $18$ \\
\hline
Total & $10$ & $22$ & $32$ \\
\hline
\end{tabular}
\end{center}
\cas{b)}
\explique{\res{$12$ élèves} ne font ni musique ni dessin.}{case « pas de musique, pas de
dessin »}
\cas{c)}
\explique{$M\cap\overline{D}$ : \res{« l'élève fait de la musique mais pas de
dessin »}.}{$\overline{D}$ : « il ne fait pas de dessin »}
\explique{$P\big(M\cap\overline{D}\big)=\dfrac{10}{32}=\res{\dfrac{5}{16}}$.}{élève choisi au
hasard : équiprobabilité}
\cas{d)}
\begin{calculs}
P(M\cup D) &= P(M)+P(D)-P(M\cap D) &&\justif{propriété 3}\\
P(M\cup D) &= \dfrac{14}{32}+\dfrac{10}{32}-\dfrac{4}{32} &&\\
P(M\cup D) &= \dfrac{20}{32} &&\\
P(M\cup D) &= \res{\dfrac{5}{8}} &&\justif{on simplifie par $4$}
\end{calculs}
\explique{Vérification : $32-12=20$ élèves font au moins l'une des deux activités.}{contraire
de « ni l'un ni l'autre »}
\end{exo}

\begin{exo}[Satisfaction des usagers]
\cas{a)}
\explique{Métro : $1\,000-400-350=250$ usagers. Satisfaits : $0{,}7\times 400=280$ pour le
bus, $0{,}8\times 350=280$ pour le tram, $0{,}9\times 250=225$ pour le métro.}{pourcentages :
chapitre 7}
\begin{center}
\renewcommand{\arraystretch}{1.5}
\begin{tabular}{|l|c|c|c|c|}
\hline
 & Bus & Tram & Métro & Total \\
\hline
Satisfait & $\res{280}$ & $280$ & $225$ & $785$ \\
\hline
Non satisfait & $120$ & $70$ & $25$ & $215$ \\
\hline
Total & $400$ & $350$ & $250$ & $1\,000$ \\
\hline
\end{tabular}
\end{center}
\cas{b)}
\explique{$P(T)=\dfrac{350}{1\,000}=\res{0{,}35}$, $P(S)=\dfrac{785}{1\,000}=\res{0{,}785}$
et $P(T\cap S)=\dfrac{280}{1\,000}=\res{0{,}28}$.}{usager choisi au hasard :
équiprobabilité}
\cas{c)}
\begin{calculs}
P(T\cup S) &= P(T)+P(S)-P(T\cap S) &&\justif{propriété 3}\\
P(T\cup S) &= 0{,}35+0{,}785-0{,}28 &&\\
P(T\cup S) &= \res{0{,}855} &&
\end{calculs}
\explique{Vérification : seuls les non-satisfaits du bus et du métro ne sont ni au tram ni
satisfaits, soit $120+25=145$, et $1-0{,}145=0{,}855$.}{contraire de $T\cup S$}
\end{exo}

\partie{Probabilités conditionnelles}

\begin{exo}[Langue vivante et demi-pension]
\cas{a)}
\explique{« Parmi les demi-pensionnaires » : on ne regarde que leur ligne, $90$ germanistes
sur $300$.}{la condition est $D$}
\explique{$P_D(A)=\dfrac{90}{300}=\res{0{,}3}$.}{$\dfrac{\Card(D\cap A)}{\Card(D)}$}
\cas{b)}
\explique{« Parmi les germanistes » : on ne regarde que leur colonne, $90$
demi-pensionnaires sur $120$.}{la condition est $A$}
\explique{$P_A(D)=\dfrac{90}{120}=\res{0{,}75}$.}{$P_D(A)\ne P_A(D)$ : l'ordre compte}
\cas{c)}
\explique{$P_E(D)=\dfrac{210}{380}\approx\res{0{,}55}$.}{colonne de l'espagnol}
\explique{\res{Parmi les élèves qui apprennent l'espagnol, environ $55\,\%$ sont
demi-pensionnaires.}}{on écrit la condition en premier : « parmi … »}
\end{exo}

\begin{exo}[De la phrase à la notation]
\cas{a)}
\explique{« $45\,\%$ des demi-pensionnaires … » : $\res{P_D(S)=0{,}45}$.}{on choisit parmi
les demi-pensionnaires : condition $D$}
\explique{« parmi les élèves de l'AS, $60\,\%$ … » : $\res{P_S(D)=0{,}6}$.}{condition $S$}
\explique{« $20\,\%$ des élèves sont des demi-pensionnaires de l'AS » :
$\res{P(D\cap S)=0{,}2}$.}{pas de condition : tous les élèves}
\cas{b)}
\explique{\res{Parmi les élèves qui ne sont pas demi-pensionnaires (les externes), $30\,\%$
font partie de l'AS.}}{condition $\overline{D}$ ; on parle de $S$}
\end{exo}

\begin{exo}[Un test de dépistage]
\cas{a)}
\explique{Malades : $0{,}02\times 20\,000=\res{400}$.}{$2\,\%$ de $20\,000$}
\explique{Tests positifs chez les malades : $0{,}9\times 400=\res{360}$.}{sensibilité :
$90\,\%$ des malades}
\cas{b)}
\explique{Non-malades : $20\,000-400=19\,600$.}{les autres personnes}
\explique{Tests négatifs : $0{,}97\times 19\,600=\res{19\,012}$ ; tests positifs :
$19\,600-19\,012=\res{588}$.}{spécificité : $97\,\%$ des non-malades}
\cas{c)}%
\begin{center}
\renewcommand{\arraystretch}{1.5}
\begin{tabular}{|l|c|c|c|}
\hline
 & Test positif & Test négatif & Total \\
\hline
Malades & $360$ & $40$ & $400$ \\
\hline
Non malades & $588$ & $19\,012$ & $19\,600$ \\
\hline
Total & $948$ & $19\,052$ & $20\,000$ \\
\hline
\end{tabular}
\end{center}
\cas{d)}
\explique{On ne regarde que la colonne des $948$ tests positifs.}{la condition est $T$}
\explique{$P_T(M)=\dfrac{360}{948}\approx\res{0{,}38}$.}{calculatrice}
\cas{e)}
\explique{Faux positifs : $P_T\big(\overline{M}\big)=\dfrac{588}{948}\approx\res{0{,}62}$.}{non
malades parmi les positifs}
\explique{Le test détecte $90\,\%$ des malades, mais environ $62\,\%$ des tests positifs
sont de faux positifs : la maladie est rare, et $3\,\%$ des très nombreux non-malades
dépassent les malades détectés. Un test positif doit être confirmé.}{$P_M(T)=0{,}9$ mais
$P_T(M)\approx 0{,}38$}
\end{exo}

\begin{exo}[Avec la formule]
\cas{a)}
\explique{Équiprobabilité sur les $32$ cartes : $P(C)=\dfrac{8}{32}=\res{\dfrac{1}{4}}$,
$P(R)=\dfrac{4}{32}=\res{\dfrac{1}{8}}$.}{$8$ cœurs, $4$ rois}
\explique{$C\cap R$ : le roi de cœur, donc $P(C\cap R)=\res{\dfrac{1}{32}}$.}{une seule carte}
\cas{b)}
\begin{calculs}
P_C(R) &= \dfrac{P(C\cap R)}{P(C)} &&\justif{propriété 6}\\
P_C(R) &= \dfrac{1}{32}\times\dfrac{4}{1} &&\justif{diviser par $\frac{1}{4}$, c'est
multiplier par $4$}\\
P_C(R) &= \dfrac{4}{32} &&\\
P_C(R) &= \res{\dfrac{1}{8}} &&
\end{calculs}
\explique{Parmi les cœurs, un sur huit est un roi.}{vérification : $1$ roi sur $8$ cœurs}
\begin{calculs}
P_R(C) &= \dfrac{P(C\cap R)}{P(R)} &&\justif{propriété 6}\\
P_R(C) &= \dfrac{1}{32}\times\dfrac{8}{1} &&\justif{diviser par $\frac{1}{8}$, c'est
multiplier par $8$}\\
P_R(C) &= \dfrac{8}{32} &&\\
P_R(C) &= \res{\dfrac{1}{4}} &&
\end{calculs}
\explique{Parmi les rois, un sur quatre est un cœur.}{vérification : $1$ cœur sur $4$ rois}
\cas{c)}
\begin{calculs}
P(C\cup R) &= P(C)+P(R)-P(C\cap R) &&\justif{propriété 3}\\
P(C\cup R) &= \dfrac{8}{32}+\dfrac{4}{32}-\dfrac{1}{32} &&\justif{même dénominateur $32$}\\
P(C\cup R) &= \res{\dfrac{11}{32}} &&
\end{calculs}
\end{exo}

\partie{Arbres pondérés}

\begin{exo}[Compétition au club de tennis]
\cas{a)}
\cotefig{\explique{$P(A)=0{,}6$, $P_A(C)=0{,}3$ et $P_{\overline{A}}(C)=0{,}7$.}{« parmi
les adultes » : sachant $A$}
\explique{On complète chaque nœud à $1$, d'où $P\big(\overline{A}\big)=0{,}4$,
$P_A\big(\overline{C}\big)=0{,}7$, $P_{\overline{A}}\big(\overline{C}\big)=0{,}3$.}{somme des
branches d'un nœud : $1$}}
{\begin{tikzpicture}[arbre]
\node {}
  child {node {$A$}
    child {node {$C$} edge from parent node[poids,above] {$0{,}3$}}
    child {node {$\overline{C}$} edge from parent node[poids,below] {$0{,}7$}}
    edge from parent node[poids,above] {$0{,}6$}}
  child {node {$\overline{A}$}
    child {node {$C$} edge from parent node[poids,above] {$0{,}7$}}
    child {node {$\overline{C}$} edge from parent node[poids,below] {$0{,}3$}}
    edge from parent node[poids,below] {$0{,}4$}};
\end{tikzpicture}}
\explique{\res{Arbre ci-contre.}}{premier niveau : adulte ou non ; second : compétition ou
non}
\cas{b)}
\begin{calculs}
P(A\cap C) &= P(A)\times P_A(C) &&\justif{chemin $A$ puis $C$}\\
P(A\cap C) &= 0{,}6\times 0{,}3 &&\\
P(A\cap C) &= \res{0{,}18} &&
\end{calculs}
\cas{c)}
\begin{calculs}
P(C) &= P(A\cap C)+P\big(\overline{A}\cap C\big) &&\justif{somme des deux chemins qui
mènent à $C$}\\
P(C) &= 0{,}18+0{,}4\times 0{,}7 &&\\
P(C) &= 0{,}18+0{,}28 &&\\
P(C) &= \res{0{,}46} &&
\end{calculs}
\cas{d)}
\begin{calculs}
P\big(\overline{A}\cap\overline{C}\big) &= 0{,}4\times 0{,}3 &&\justif{chemin
$\overline{A}$ puis $\overline{C}$}\\
P\big(\overline{A}\cap\overline{C}\big) &= \res{0{,}12} &&
\end{calculs}
\end{exo}

\begin{exo}[Deux chocolats]
\cas{a)}
\cotefig{\explique{Au premier choix : $3$ noirs ($N$) et $4$ au lait ($L$) sur $7$.}{tous
les chocolats ont la même chance}
\explique{Après un noir, il reste $2$ noirs et $4$ au lait ; après un chocolat au lait,
$3$ noirs et $3$ au lait.}{on a mangé le premier : il en reste $6$}}
{\begin{tikzpicture}[arbre]
\node {}
  child {node {$N$}
    child {node {$N$} edge from parent node[poids,above] {$\frac{2}{6}$}}
    child {node {$L$} edge from parent node[poids,below] {$\frac{4}{6}$}}
    edge from parent node[poids,above] {$\frac{3}{7}$}}
  child {node {$L$}
    child {node {$N$} edge from parent node[poids,above] {$\frac{3}{6}$}}
    child {node {$L$} edge from parent node[poids,below] {$\frac{3}{6}$}}
    edge from parent node[poids,below] {$\frac{4}{7}$}};
\end{tikzpicture}}
\explique{\res{Arbre ci-contre.}}{chaque nœud a pour somme $1$}
\cas{b)}
\begin{calculs}
P(\text{deux noirs}) &= \dfrac{3}{7}\times\dfrac{2}{6} &&\justif{chemin $N$ puis $N$}\\
P(\text{deux noirs}) &= \dfrac{6}{42} &&\\
P(\text{deux noirs}) &= \res{\dfrac{1}{7}} &&\justif{on simplifie par $6$}
\end{calculs}
\begin{calculs}
P(\text{deux au lait}) &= \dfrac{4}{7}\times\dfrac{3}{6} &&\justif{chemin $L$ puis $L$}\\
P(\text{deux au lait}) &= \dfrac{12}{42} &&\\
P(\text{deux au lait}) &= \res{\dfrac{2}{7}} &&\justif{on simplifie par $6$}
\end{calculs}
\cas{c)}
\begin{calculs}
P(\text{même sorte}) &= \dfrac{1}{7}+\dfrac{2}{7} &&\justif{somme des deux chemins}\\
P(\text{même sorte}) &= \res{\dfrac{3}{7}} &&
\end{calculs}
\cas{d)}
\explique{Le contraire de « au moins un noir » est « deux au lait ».}{aucun noir}
\begin{calculs}
P(\text{au moins un noir}) &= 1-\dfrac{2}{7} &&\justif{probabilité du contraire}\\
P(\text{au moins un noir}) &= \res{\dfrac{5}{7}} &&
\end{calculs}
\end{exo}

\begin{exo}[Deux tickets à gratter]
\cas{a)}
\cotefig{\explique{Chaque ticket est gagnant avec la probabilité $\frac{1}{4}$, perdant avec
$\frac{3}{4}$, quel que soit le résultat de l'autre.}{les mêmes probabilités au second
niveau}
\explique{\res{Arbre ci-contre.}}{$G$ puis $G$ : deux tickets gagnants}}
{\begin{tikzpicture}[arbre]
\node {}
  child {node {$G$}
    child {node {$G$} edge from parent node[poids,above] {$\frac{1}{4}$}}
    child {node {$\overline{G}$} edge from parent node[poids,below] {$\frac{3}{4}$}}
    edge from parent node[poids,above] {$\frac{1}{4}$}}
  child {node {$\overline{G}$}
    child {node {$G$} edge from parent node[poids,above] {$\frac{1}{4}$}}
    child {node {$\overline{G}$} edge from parent node[poids,below] {$\frac{3}{4}$}}
    edge from parent node[poids,below] {$\frac{3}{4}$}};
\end{tikzpicture}}
\cas{b)}
\explique{$P(\text{deux gagnants})=\dfrac{1}{4}\times\dfrac{1}{4}=\res{\dfrac{1}{16}}$.}{chemin
$G$ puis $G$}
\cas{c)}
\explique{Un seul gagnant : chemins $G$ puis $\overline{G}$, et $\overline{G}$ puis $G$.}{deux
chemins}
\begin{calculs}
P(\text{un seul}) &= \dfrac{1}{4}\times\dfrac{3}{4}+\dfrac{3}{4}\times\dfrac{1}{4}
  &&\justif{somme des deux chemins}\\
P(\text{un seul}) &= \dfrac{3}{16}+\dfrac{3}{16} &&\\
P(\text{un seul}) &= \dfrac{6}{16} &&\\
P(\text{un seul}) &= \res{\dfrac{3}{8}} &&\justif{on simplifie par $2$}
\end{calculs}
\cas{d)}
\explique{Le contraire de « au moins un gagnant » est « aucun gagnant » :
$\dfrac{3}{4}\times\dfrac{3}{4}=\dfrac{9}{16}$.}{chemin $\overline{G}$ puis $\overline{G}$}
\begin{calculs}
P(\text{au moins un}) &= 1-\dfrac{9}{16} &&\justif{probabilité du contraire}\\
P(\text{au moins un}) &= \res{\dfrac{7}{16}} &&
\end{calculs}
\explique{Vérification : $\dfrac{1}{16}+\dfrac{6}{16}=\dfrac{7}{16}$.}{deux gagnants ou un seul}
\end{exo}

\begin{exo}[Un test pour les chats]
\cas{a)}
\cotefig{\explique{$P(M)=0{,}01$ ; sensibilité $P_M(T)=0{,}98$ ; spécificité
$P_{\overline{M}}\big(\overline{T}\big)=0{,}95$.}{on traduit l'énoncé}
\explique{On complète : $P\big(\overline{M}\big)=0{,}99$,
$P_M\big(\overline{T}\big)=0{,}02$ et $P_{\overline{M}}(T)=0{,}05$.}{somme des branches d'un
nœud : $1$}
\explique{\res{Arbre ci-contre.}}{$0{,}05$ : les faux positifs}}
{\begin{tikzpicture}[arbre]
\node {}
  child {node {$M$}
    child {node {$T$} edge from parent node[poids,above] {$0{,}98$}}
    child {node {$\overline{T}$} edge from parent node[poids,below] {$0{,}02$}}
    edge from parent node[poids,above] {$0{,}01$}}
  child {node {$\overline{M}$}
    child {node {$T$} edge from parent node[poids,above] {$0{,}05$}}
    child {node {$\overline{T}$} edge from parent node[poids,below] {$0{,}95$}}
    edge from parent node[poids,below] {$0{,}99$}};
\end{tikzpicture}}
\cas{b)}
\begin{calculs}
P(M\cap T) &= 0{,}01\times 0{,}98 &&\justif{chemin $M$ puis $T$}\\
P(M\cap T) &= \res{0{,}0098} &&
\end{calculs}
\begin{calculs}
P(T) &= P(M\cap T)+P\big(\overline{M}\cap T\big) &&\justif{somme des deux chemins}\\
P(T) &= 0{,}0098+0{,}99\times 0{,}05 &&\\
P(T) &= 0{,}0098+0{,}0495 &&\\
P(T) &= \res{0{,}0593} &&
\end{calculs}
\cas{c)}
\begin{calculs}
P_T(M) &= \dfrac{P(M\cap T)}{P(T)} &&\justif{propriété 6}\\
P_T(M) &= \dfrac{0{,}0098}{0{,}0593} &&\\
P_T(M) &\approx \res{0{,}17} &&\justif{calculatrice}
\end{calculs}
\cas{d)}
\explique{$P_M(T)=0{,}98$ mais $P_T(M)\approx 0{,}17$ : \res{un test positif ne suffit pas},
car environ $83\,\%$ des chats testés positifs ne sont pas malades.}{la maladie est rare :
les faux positifs sont plus nombreux que les vrais}
\end{exo}

\partie{Loi des grands nombres et simulation}

\begin{exo}[Lire une simulation]
\cas{a)}
\explique{Une face gagnante sur $4$, dé équilibré : $p=\res{\dfrac{1}{4}}$, soit
$0{,}25$.}{équiprobabilité}
\cas{b)}
\explique{\texttt{gains} compte les lancers où le dé donne $4$ ; la fonction renvoie
\res{\texttt{gains / n}, la fréquence observée des gains sur $n$ lancers}.}{nombre de gains
divisé par $n$}
\cas{c)}
\explique{Plus $n$ est grand, plus la fréquence est proche de $0{,}25$ ; avec $n=10$, elle
en est loin ($0{,}4$). C'est \res{la loi des grands nombres}.}{propriété 8}
\cas{d)}
\explique{Pour $n=100$ : l'écart vaut $0{,}25-0{,}21$, soit $0{,}04$, et
$\dfrac{1}{\sqrt{100}}=0{,}1$ : \res{l'écart est inférieur à $0{,}1$}.}{$\sqrt{100}=10$}
\explique{Pour $n=1\,000$ : l'écart vaut $0{,}262-0{,}25$, soit $0{,}012$, et
$\dfrac{1}{\sqrt{1\,000}}\approx 0{,}032$ : \res{l'écart est inférieur à
$0{,}032$}.}{calculatrice pour $\sqrt{1\,000}$}
\end{exo}

\begin{exo}[Simuler un contrôle de qualité]
\cas{a)}
\begin{python}
from random import random

def livre():
    if random() < 0.04:
        return 1
    else:
        return 0

def frequence(n):
    defauts = 0
    for k in range(n):
        defauts = defauts + livre()
    return defauts / n
\end{python}
\explique{\res{\texttt{random() < 0.04}, \texttt{return 0} et \texttt{return defauts / n}}.}{\texttt{random()
< 0.04} est vrai avec la probabilité $0{,}04$}
\cas{b)}
\explique{D'après la loi des grands nombres, \texttt{frequence(10000)} sera proche de
\res{$0{,}04$}.}{$4\,\%$ des livres ont un défaut}
\cas{c)}
\explique{$\dfrac{1}{\sqrt{10\,000}}=\dfrac{1}{100}=0{,}01$ : dans environ $95\,\%$ des cas,
la fréquence est entre $0{,}04-0{,}01$ et $0{,}04+0{,}01$, soit \res{entre $0{,}03$ et
$0{,}05$}.}{propriété 8 ; $\sqrt{10\,000}=100$}
\end{exo}

\begin{exo}[Estimer une probabilité]
\cas{a)}
\explique{$f=\dfrac{1\,550}{2\,500}=\res{0{,}62}$.}{nombre de « tête » divisé par le nombre
de lancers}
\cas{b)}
\explique{$\dfrac{1}{\sqrt{2\,500}}=\dfrac{1}{50}=\res{0{,}02}$.}{$50\times 50=2\,500$}
\explique{Dans environ $95\,\%$ des cas, l'écart entre $f$ et $p$ est inférieur à $0{,}02$ :
$p$ est vraisemblablement \res{entre $0{,}60$ et $0{,}64$}.}{$0{,}62-0{,}02$ et
$0{,}62+0{,}02$}
\cas{c)}
\explique{Il faut $\dfrac{1}{\sqrt{n}}=0{,}01$, soit $\sqrt{n}=100$, donc
$n=\res{10\,000}$ lancers.}{diviser l'erreur par $2$ demande $4$ fois plus de lancers}
\end{exo}

\end{document}
